\section{由共享标签的价格推导下界}\label{sec:boundary}

考虑两个 \(n\) 阶方阵的乘法，把两个输入和输出分别放在三条臂长均为一的三脚架的三个端点。三个矩阵指标各由相距 \(2\) 的两个端点共享，因此乘法张量的价格为 \(6\log n\)。一个秩为 \(r\) 的分解可由三个端点共享的、具有 \(r\) 个取值的标签提供这个张量。若此共享标签的费率为 \(8/3\)，则其价格为 \((8/3)\log r\)，由价格的单调性可得 \(r\ge n^{9/4}\)。我们将在第~\ref{sec:motion}~节构造的族中找到这样的三脚架。先考察一条臂远长于另两条臂时的共享标签费率。以下原理把所得等式推广到三臂等长的情形。

\begin{proposition}[三脚架原理]\label{prop:tripod}
在代数闭域上放置三个站点，使对每一对 \(i\ne j\)，站点间距为 \(a_i+a_j\)，其中 \(a_1,a_2,a_3\ge0\)。其共享标签费率 \(D(a)\) 满足
\[
       D(a)\ \ge\ D_*(a)
       :=a_1+a_2+a_3-\bigl(a_1^{-1}+a_2^{-1}+a_3^{-1}\bigr)^{-1}.
\]
只要有一个 \(a_i=0\)，右侧的倒数项就取零。

现固定谱空间中的三条无限射线，它们构成一个以 \(O\) 为原点的度量三脚架，并假设其价格在径向伸缩下齐次。在 \(D(a)\) 中，把站点 \(i\) 放在射线 \(i\) 上、距 \(O\) 为 \(a_i\) 的位置。若 \(D(a)-D_*(a)\) 在某个正三元组处为零，则它在每个正三元组处都为零。若固定其中两条臂的正长度而让第三条臂趋于无穷时，该差值趋于零，则结论也成立。
\end{proposition}

第一个断言适用于满足指定距离的任意配置。第二个断言比较同一个三脚架上的不同配置。第~\ref{sec:profiles}~至~\ref{sec:contact}~节通过研究多项式乘法的价格如何随次数和系数的位置变化，证明这两个断言。

\subsection{矩形矩阵乘法的解码器}

在第~\ref{sec:motion}~节的代数闭域 \(F\) 上，将一个 \(s\times q\) 矩阵与一个 \(q\times s\) 矩阵相乘。取四次幂 \(N\)，令 \(s=N^{1/2}\)、\(q=N^{1/4}\)。乘积有 \(N\) 个元素。将输出放在 \(A\)，两个输入放在 \(C,D'\)，对应的臂长依次为 \(c,a,b\)。外部指标在 \(AC,AD'\) 上给出大小为 \(s\) 的连线，内部指标在 \(CD'\) 上给出大小为 \(q\) 的连线。因此，将乘积的价格除以 \(\log N\)，所得数满足
\[
   c+\tfrac34(a+b)
   \ \le\ c+a+b-\frac{ab}{a+b}
   \ \le\ D_*(c,a,b)\ \le\ D(c,a,b).
\]
在 \(a=b=0\) 时约定 \(ab/(a+b)=0\)。这就给出了与半宽矩阵乘法对应的矩形界 \cite{uznanski,abvz}。

特别地，对每个 \(\eta>0\)，都有任意大的这种 \(N\)，使精确张量秩至多为 \(N^{1+\eta}\)。事实上，从充分大的四次幂 \(N_0\) 和一个大小为 \(\lceil N_0^{1+\eta/2}\rceil\) 的全对角张量出发。在每个谱点，上述不等式都给出与该对角张量的严格谱比较。乘法张量的每个张量幂都是对应对角张量幂的有界多个副本的限制。幂次充分大时，指数 \(1+\eta\) 足以吸收副本数和取整的影响。

\begin{lemma}[具有短行的解码器]\label{lem:secondary}
对每个 \(\epsilon>0\)，都存在 \(\delta>0\)，使得下述结论对任意大的某些四次幂 \(N\) 成立：这一乘法的 \(N\) 个输出双线性型 \(B_e\) 均有如下分解：
\[
    B_e=S_e+R_e,\qquad
    \operatorname{rank}S_e\le N^\epsilon,\qquad
    \dim\langle R_e:1\le e\le N\rangle\le N^{1-\delta}.
\]
\end{lemma}

这两个界各有用途。共享输出指标 \(e\) 后，每个 \(S_e\) 都可用少量简单乘积的和计算。修正项 \(R_e\) 单独计算可能很昂贵，但它们张成的整个空间可以沿一条连线传输。为得到对基本规格一致的界，我们用一个工作寄存器实现解码器的张量幂。克罗内克积的相关非刚性估计见 \cite{almanrigidity,kivva}。

\begin{proof}
只需考虑 \(0<\epsilon<1\)。固定整数 \(K>10/\epsilon\)。设基本规格有 \(d\) 个输出元素，秩满足 \(d\le r\le d^{1+\eta}\)。令 \(M\) 为把输出表示成 \(r\) 个简单双线性乘积的线性组合的 \(d\times r\) 矩阵。稍后选取 \(\eta\)。第 \(Km\) 个张量幂有 \(N=d^{Km}\) 个输出元素。为每个解码块 \(M^{\otimes m}\) 补上零行和零列，使其成为边长为 \(2^t\) 的方阵，其中 \(t=\lceil m\log_2 r\rceil\)。补零后的整个解码器的边长为 \(T=2^{Kt}\)。每个指标由 \(K\) 个寄存器组成，每个寄存器有 \(t\) 比特。

将第一个块矩阵作用于第一个寄存器。对之后的每一个块，先把其寄存器与第一个交换，在第一个寄存器上施加块矩阵，再交换回来。一次寄存器交换同时交换 \(t\) 对比特。一次比特交换置换四个字母中的两个，即 \(01\) 和 \(10\)。在任意域上，它的矩阵都可分解为固定数目的初等剪切和对角符号变换。因此一次寄存器交换由固定数目的阶段组成，每一阶段在所有比特位置施加同一局部剪切或对角符号变换。

对于剪切，把非对角元素指向的字母称为高位字母。固定一个充分小的 \(\rho>0\)，在该阶段的前后分别作投影，仅保留寄存器对中高位字母的个数介于 \(t/4-\rho t\) 和 \(t/4+\rho t\) 之间的坐标。由二项分布的尾界，当 \(t\) 充分大时，被丢弃的坐标比例至多为 \(e^{-c_\rho t}\)。将 \(O(K)\) 次投影的差逐项作望远镜求和，得到解码器被丢弃部分的秩至多为 \(O(K)Te^{-c_\rho t}\le T^{1-\gamma}\)，其中 \(\gamma=\gamma(K,\rho)>0\)。在剪切阶段被保留的路径上，每一次改变都产生一个高位字母，因此至多有 \(2\rho t\) 个位置发生改变。

块矩阵只作用于第一个寄存器。其他每个寄存器恰参与两次交换，所以在保留的路径上，其输入与输出间的汉明距离至多为 \(C\rho t\)，其中 \(C\) 为绝对常数。因此，保留的解码器的各输出行中，非零元素个数至多为
\[
       2^t\left(\sum_{j\le C\rho t}\binom{t}{j}\right)^{K-1}
       \ \le\ T^{1/K+o_{\rho\downarrow0}(1)}
\]
选取 \(\rho\) 使该指数小于 \(\epsilon/2\)，再取充分小的 \(\eta\)。由于 \(T\le 2^K N^{1+\eta}\)，对一切充分大的 \(m\)，上述两个界分别成为每个保留行至多有 \(N^\epsilon\) 个元素、被丢弃矩阵的秩至多为 \(N^{1-\delta}\)，其中 \(\delta=\gamma/2\)。去掉补上的零行零列，把两个矩阵分别作用于简单双线性乘积的列表。保留下来的行给出双线性型 \(S_e\)，而被丢弃的行空间的像张成 \(R_e\)。
\end{proof}

\subsection{碰撞附近的两个输出副本}

取第~\ref{sec:motion}~节中站点 \(4\) 趋近站点 \(3\) 的碰撞，把站点 \(1,2,3,4\) 分别记为 \(A,B,C,D'\)。在参数 \(r\) 处，记 \(L=2\ell\)、\(h=L-r\)，其中 \(0<h<L\)。\(AB\) 与 \(CD'\) 的内部距离均为 \(h\)，四个跨对距离均为 \(L\)。用 \(D_\ell(h)\) 表示四站点的共享标签费率。碰撞时 \(D_\ell(0)=L\)，而在一般的四站点配置下 \(D_\ell(L)=D_\ell\)。由凸性得到
\[
                   D_\ell(h)\le L+\frac hL(D_\ell-L).
\]

在 \(A\) 和 \(B\) 各放置矩形乘法的输出指标 \(e=(i,i')\) 的一个副本，两个输入仍放在 \(C,D'\)。所得张量为
\[
 U_N=\sum_e e_e\otimes e_e\otimes B_e
     =\sum_{\substack{1\le i,i'\le s\\1\le j\le q}}
        e_{ii'}\otimes e_{ii'}\otimes e_{ij}\otimes e_{ji'} .
\]
指标 \(i\) 由 \(A,B,C\) 共享，指标 \(i'\) 由 \(A,B,D'\) 共享，而 \(j\) 沿短连线 \(CD'\) 传递。以 \(D_{ABC},D_{ABD'}\) 表示两个三元组的费率，则
\[
   \frac{H(U_N)}{\log N}=\frac{D_{ABC}+D_{ABD'}}2+\frac h4 .
\]
两个三元组的臂长均为 \((h/2,h/2,L-h/2)\)。命题~\ref{prop:tripod} 中的倒数修正项至多为 \(h/4\)，故两者的费率均至少为 \(L+h/4\)。

采用引理~\ref{lem:secondary} 的分解 \(B_e=S_e+R_e\)。以 \(\log N\) 为单位，诸 \(S_e\) 形成的张量的成本至多为 \(D_\ell(h)+\epsilon h\)。先在四个站点共享标签 \(e\)，然后在 \(CD'\) 上用一条大小至多为 \(N^\epsilon\) 的连线提供指定 \(e\) 所需的简单乘积。诸 \(R_e\) 形成的张量沿切分 \(AB\mid CD'\) 的矩阵秩至多为 \(N^{1-\delta}\)。用连接 \(A\) 和 \(C\) 的一条连线提供这一秩，再由两对站点各自内部的连线把维数至多为 \(N\) 的张量腿分别送到 \(B\) 和 \(D'\)。这样做的成本至多为 \((1-\delta)L+2h\)。

取 \(h/L<\delta/2\)。此时修正项的成本小于 \(L\)，由限制到一对分属两组的站点可得 \(D_\ell(h)\ge L\)。在固定配置下，对普通和使用价格界只需付出与 \(N\) 无关的加法常数。令 \(N\) 沿引理~\ref{lem:secondary} 的规格趋于无穷，再在固定 \(h/L\) 时应用式~\eqref{eq:motion-upper}，可得
\[
  \frac{D_{ABC}+D_{ABD'}}2+\frac h4
  \ \le\ D_\ell(h)+\epsilon h
  \ \le\ L+\bigl(\tfrac12+\epsilon+o_{\ell\downarrow0}(1)\bigr)h .
\]
减去下界 \(D_{ABD'}\ge L+h/4\) 得
\[
        D_{ABC}\le L+h/4+\bigl(2\epsilon+o_{\ell\downarrow0}(1)\bigr)h.
\]

前三个站点位于从深度为零的点出发的固定射线 \(z_1,z_2,z_3\) 上。为在该点取切三脚架，把站点放在深度 \(t a_i\)，将价格除以 \(t\)，并在谱超极限中令 \(t\downarrow0\)。在这些固定射线上，每个非零张量经缩放的价格都有极限，因为原价格关于 \(t\) 为凸函数且在零处消失。极限价格在径向伸缩下齐次。

令 \(t=h/2\) 并固定 \(b=2L/h-1\)。三个深度变成 \(t(1,1,b)\)。用 \(t\) 除上一个不等式得到 \(D(1,1,b)\le b+3/2+4\epsilon\)，而
\[
             D_*(1,1,b)=b+\tfrac32+\frac{1}{4b+2}.
\]
因此，对一切充分大的有理数 \(b\) 都有 \(0\le D(1,1,b)-D_*(1,1,b)\le4\epsilon\)。由连续性，该式对充分大的实数 \(b\) 也成立。先令 \(b\) 趋于无穷，再令 \(\epsilon\) 趋于零，就得到在这一固定三脚架上无穷远处的接触。命题~\ref{prop:tripod} 现给出 \(D(1,1,1)=3-1/3=8/3\)。本节开头的三连线价格比较随即给出 \(r\ge n^{9/4}\)，从而证明了定理~\ref{thm:motion}。
